\documentclass[10pt,a4paper]{amsart}
\usepackage[margin=19mm]{geometry}
\usepackage{amsmath,amssymb,mathtools}
\usepackage[scr=rsfs,cal=euler]{mathalfa}
\usepackage{xfrac}
\usepackage[unicode,hidelinks,bookmarks=false]{hyperref}
\newcommand{\ie}{i.e.\ }
\newcommand{\cf}{cf.\ }
\newcommand{\resp}{respectively\ }
\newcommand{\Subsection}{\subsection}
\providecommand{\nobreakdash}{\nobreak\hbox{-}}
\setlength{\emergencystretch}{2em}
\multlinegap0pt
\usepackage{bm}
\usepackage{bbm}
\usepackage{dsfont}
\usepackage{xspace}
\usepackage{cancel}
\usepackage{mathtools}
\usepackage{amsthm}
\makeatletter
\@ifundefined{theorem}{\newtheorem{theorem}{Theorem}}{}
\@ifundefined{assumption}{\newtheorem{assumption}[theorem]{Assumption}}{}
\@ifundefined{lemma}{\newtheorem{lemma}[theorem]{Lemma}}{}
\makeatother
\newcommand{\Asf}{\mathsf A}
\newcommand{\Dom}{\operatorname{Dom}}
\newcommand{\FC}{\mathcal{FC}_b^\infty}
\newcommand{\Hh}{\mathcal H}
\newcommand{\Ran}{\operatorname{Ran}}
\title[Noncommutative Anisotropic Diffusion in Hilbert Space. I. Th]{Noncommutative Anisotropic Diffusion in Hilbert Space. I. The Consistent A-Geometry, Mosco Stability, and the Weak Bridge}
\author{E. Yu. Shchetinin, A. A. Shevchuk, S. I. Salpagarov}
\date{Source: arXiv:2606.28964v1 (CC BY 4.0)}
\begin{document}
\maketitle

\noindent\emph{Source and licence.} Noncommutative Anisotropic Diffusion in Hilbert Space. I. The Consistent A-Geometry, Mosco Stability, and the Weak Bridge. Version arXiv:2606.28964v1 (CC BY 4.0), distributed under the Creative Commons Attribution 4.0 International licence, as recorded in the arXiv licence metadata for this exact version and shown on the abstract page https://arxiv.org/abs/2606.28964v1. Full rights notice: licenses/arxiv-2606.28964-CC-BY-4.0.txt.\par\smallskip

\noindent\textit{Editorial scope.} Counted unit: the Theorem whose source label is \texttt{thm:bridge} in the article (source0.tex lines 609--625, source label \texttt{thm:bridge}), delivered together with the article's own complete original proof of it (source0.tex lines 627--672, one \texttt{proof} environment of 46 lines). No mathematics is written, repaired or re-derived: statement and proof are the article's own text line for line. Required context retained uncounted: lem:dual (lines 470--481); ass:bridge (lines 599--607). Every definition, display and label of those lines is retained unchanged. Omitted from this package and not redistributed: 1--469, 482--598, 608--608, 626--626, 673--1722. Nothing inside a retained excerpt is omitted or altered: every retained body is delivered exactly as the article's lines are written, so no omission receipt of any kind accompanies this package. Citations: the retained excerpts carry no citation command at all, so no bibliography entry of the article is transcribed and no reference is redistributed. Reference adaptation: none was needed and none was made; every reference inside the delivered unit resolves inside it, so no unresolved reference or citation is produced. Printed slips kept literal: no printed word, symbol, hypothesis or number of the counted statement or of its proof was altered. Layout and class: the article's own class is \texttt{article} with packages including arxiv, inputenc, fontenc, hyperref, url, amsmath, amssymb, amsthm, amsfonts, nicefrac, microtype, graphicx; the wrapper is the lane's pinned \texttt{amsart} editorial wrapper at 10pt on A4, which restates the theorem-like environments the retained text needs and adds the article's own macro definitions that the retained text needs (\texttt{\textbackslash Asf}, \texttt{\textbackslash Dom}, \texttt{\textbackslash FC}, \texttt{\textbackslash Hh}, \texttt{\textbackslash Ran}), with extra wrapper packages (bm, bbm, dsfont, xspace, cancel, mathtools). The delivered numbering is the unit's own. Assets: no retained span contains a figure environment, a graphics-inclusion command, a PDF-page inclusion, a table or a TikZ picture; no image file is copied into this package, the package declares no asset dependency and no image file is redistributed with it. Licence: version arXiv:2606.28964v1 of this article is distributed under the Creative Commons Attribution 4.0 International licence, as recorded in the arXiv licence metadata for this exact version and shown on the abstract page https://arxiv.org/abs/2606.28964v1. A full rights notice is delivered at licenses/arxiv-2606.28964-CC-BY-4.0.txt. Characters: every retained excerpt is pure ASCII, so the delivered engine, XeLaTeX with its default Latin Modern text font, reports no missing character.
\begin{lemma}[Integral duality]
\label{lem:dual}
Let \(z\in L^2(\nu;\Hh)\) with
\(z(x)\in\Ran \Asf(x)^{1/2}\) for \(\nu\)-a.e.\ \(x\). Then
\[
        \sup_{\phi\in\FC}
        \frac{\int\langle z(x),\nabla\phi(x)\rangle\,d\nu(x)}
        {\mathcal E_{A,\nu}(\phi,\phi)^{1/2}}
        =
        \|z\|_{\Asf^{-1},L^2(\nu)} .
\]
\end{lemma}

\begin{assumption}[The right-hand-side space]
\label{ass:bridge}
For almost every \(t\) the functional \(F_t\) is continuous on
\(\Dom(\mathcal E_{A,\rho_t})/\mathbb R\):
\[
        |\langle F_t,\phi\rangle|
        \le C_t\mathcal E_{A,\rho_t}(\phi,\phi)^{1/2}.
\]
\end{assumption}

\begin{theorem}[The weak bridge]
\label{thm:bridge}
Under Assumption~\ref{ass:bridge} there exists
\(\Phi_t\in\Dom(\mathcal E_{A,\rho_t})/\mathbb R\) such that
\[
        \mathcal E_{A,\rho_t}(\Phi_t,\phi)=\langle F_t,\phi\rangle
\]
for all \(\phi\). The field
\[
        v_t(x)=\Asf(x)\nabla\Phi_t(x)
\]
realizes the backward weak form and satisfies
\[
        \int\|v_t(x)\|_{\Asf(x)^{-1}}^2\,d\nu_t(x)
        \le C_t^2.
\]
\end{theorem}

\begin{proof}
Let
\[
        \mathcal V_t=\overline{\FC/\mathbb R}^{\,\mathcal E_{A,\rho_t}}
\]
be the quotient space modulo constants, equipped with the norm
\(\|\phi\|_{\mathcal V_t}=\mathcal E_{A,\rho_t}(\phi,\phi)^{1/2}\). On
\(\mathcal V_t\) the bilinear form
\[
        B_t(\Phi,\phi)=\mathcal E_{A,\rho_t}(\Phi,\phi)
\]
is continuous and coercive:
\[
        |B_t(\Phi,\phi)|\le \|\Phi\|_{\mathcal V_t}\|\phi\|_{\mathcal V_t},
        \qquad
        B_t(\Phi,\Phi)=\|\Phi\|_{\mathcal V_t}^2 .
\]
Assumption~\ref{ass:bridge} means that \(F_t\in\mathcal V_t'\) and
\(\|F_t\|_{\mathcal V_t'}\le C_t\). By the Lax--Milgram theorem there exists a
unique \(\Phi_t\in\mathcal V_t\) such that
\[
        \mathcal E_{A,\rho_t}(\Phi_t,\phi)=\langle F_t,
\phi\rangle
        \quad\forall\phi\in\mathcal V_t,
\]
and \(\|\Phi_t\|_{\mathcal V_t}\le C_t\). Define the field in the negative
energy space by
\[
        v_t(x)=\Asf(x)\nabla\Phi_t(x).
\]
Then for cylindrical \(\phi\)
\[
        \int\langle v_t,\nabla\phi\rangle\,d\nu_t
        =\mathcal E_{A,\rho_t}(\Phi_t,
\phi)=\langle F_t,
\phi\rangle,
\]
that is, the field realizes the backward weak form. Since \(v_t=\Asf(x)\nabla\Phi_t\), the duality of
Lemma~\ref{lem:dual} gives
\[
        \int\|v_t(x)\|_{\Asf(x)^{-1}}^2d\nu_t(x)
        =\mathcal E_{A,\rho_t}(\Phi_t,
\Phi_t)
        \le C_t^2.
\]
\end{proof}

\end{document}
